\section{Derivación de los canales de incidencia \texorpdfstring{\(90\)}{90} y \texorpdfstring{\(120\)}{120} a partir de las semicoronas tríticas del Topological Phase Kernel}
\label{sec:sucesor-102-incidencia-tpk-90-120}

Los cardinales \(90\) y \(120\) se generan en la dinámica interna del
Topological Phase Kernel antes de toda evaluación exponencial, realización
hexada--octada o lectura metrológica. Su procedencia es la partición orientada
de las \(9\) fases y la incidencia de sus pares activos; los canales
espectrales \(q_\pm\) actuarán posteriormente sobre estos objetos ya
construidos.

\subsection{Partición trítica de la fase nonádica en semicoronas orientadas}

En la carta canónica \(D_9=\{1,\ldots,9\}\), el selector trítico de fase
determina la descomposición disjunta
\[
 H_+=\{1,4,7\},\qquad
 H_-=\{2,5,8\},\qquad
 H_0=\{3,6,9\}.
\]
Las clases de \(H_+\) activan la hoja aditiva, las de \(H_-\) la hoja
multiplicativa y las de \(H_0\) conservan el umbral. Definimos
\[
 H_{\mathrm{act}}=H_+\sqcup H_-,
 \qquad |H_{\mathrm{act}}|=6,
 \qquad \bigl|P_2(H_{\mathrm{act}})\bigr|=\binom62=15,
\]
donde \(P_2(H_{\mathrm{act}})\) denota el conjunto de pares no ordenados de
fases activas. Esta partición conserva hoja, orientación y origen de fase; no
es una clasificación retrospectiva inducida por las realizaciones físicas.

\subsection{Canal de incidencia de cardinal \texorpdfstring{\(90\)}{90}}

La incidencia interna entre una fase activa distinguida y un par activo es
\[
 \mathcal F_{90}
 :=H_{\mathrm{act}}\times P_2(H_{\mathrm{act}}).
\]
Por construcción,
\[
 |\mathcal F_{90}|=6\cdot15=90.
\]
El cardinal \(90\) procede, por tanto, de la semicorona trítica del TPK. No
es una entrada de Barbero--Immirzi, de la ley de masas ni de una tabla
metrológica.

\subsection{Canal de incidencia de cardinal \texorpdfstring{\(120\)}{120}}

Al extender el primer factor a las \(8\) posiciones que preceden al retorno
nonádico se obtiene
\[
 O_{\mathrm{ph}}=\{1,\ldots,8\},
 \qquad
 \mathcal F_{120}
 :=O_{\mathrm{ph}}\times P_2(H_{\mathrm{act}}),
\]
y, en consecuencia,
\[
 |\mathcal F_{120}|=8\cdot15=120.
\]
La diferencia entre \(\mathcal F_{90}\) y \(\mathcal F_{120}\) conserva la
información de frontera: el segundo canal incorpora las \(8\) posiciones
anteriores al retorno, mientras que el primero conserva únicamente las
\(6\) fases operativas.

\subsection{Compatibilidad dodecafásica y antecedencia causal respecto de la evaluación espectral}

La rueda dodecafásica empareja las \(2\) semicoronas orientadas y registra la
descomposición \(1+5+6=12\), invertible sobre su imagen. Así, los cardinales
de incidencia se fijan en el orden causal
\[
 \mathrm{APP}\longrightarrow\mathrm{TRIT}\longrightarrow\mathrm{TPK}
 \longrightarrow(H_+,H_-,H_0)
 \longrightarrow(\mathcal F_{90},\mathcal F_{120}).
\]
Sólo después, una vez construidas las coordenadas angulares conjugadas, los
canales \(q_\pm\) evalúan espectralmente estas incidencias mediante
\[
 V_{90,120}
 =(I-T^{90})(I-T^{120})^{-1}
 =r_+P_++r_-P_-,
 \qquad
 r_\pm=\frac{1-q_\pm^{90}}{1-q_\pm^{120}}.
\]
La evaluación \(V_{90,120}\) no genera los cardinales que aparecen en sus
exponentes: publica operatoriamente los canales de incidencia producidos con
anterioridad por el TPK.
